% Appendix H: public hyperspectral replication of Experiments 3 and 4 (Pavia University).
% Design written 2026-09-11 from code/experiments/hsi_data.py, the EXP3_DATASET / --dataset
% switches of the Exp 3 / Exp 4 code, and REFOCUS_PLAN_2026-09-10.md section 9. Results
% paragraphs are appended as the registered runs complete.

\subsection{Purpose and registration}
Our tissue spectra are not public. This appendix repeats Experiment~3 (constructed
context: the mechanism) and Experiment~4 (natural context: usage) on the ROSIS Pavia
University scene, a standard public hyperspectral benchmark ($610\times340$ pixels,
$103$ bands, nine labelled classes, $42{,}776$ labelled pixels), with the same code,
protocols, arms and pre-registered predictions, so that every number can be reproduced
from the released scripts. The scene, the spatial split, the class pair and the
disposition rules were fixed before the first registered run (pre-registration
section~9); the readiness scan and the per-pixel sweep are descriptive stages that
precede the registered runs, as on tissue.

\subsection{Spatial split, sampling and preprocessing}
The patient-level split of the tissue experiments has no counterpart in a single scene,
so we cut the scene into $20\times20$-pixel tiles and assign tiles at random (seed $0$)
to training, validation and test in the proportion $0.6/0.2/0.2$ ($208/70/70$ tiles
carrying labels), requiring at least $60$ admissible pixels per class in each split; a
pixel is admissible only if its whole $3\times3$ neighbourhood lies inside its tile, so
no evaluation neighbourhood overlaps a training one ($34{,}386$ admissible pixels). Per
tile and class up to $400/200/200$ pixels are sampled with a fixed seed: $19{,}870$
training, $6{,}172$ validation and $6{,}407$ test pixels. Residual spatial
autocorrelation across tile borders is a limitation shared by every single-scene
benchmark and is not removed here. Reflectance counts are cast to float; per-feature
standardization and PCA are fitted on the training split by the experiment code,
exactly as for the tissue spectra; there is no per-pixel normalization, denoising or
derivative.

\subsection{Experiment 3 on Pavia: constructed context}
The pair is Meadows ($+1$) against Trees ($-1$), the two largest vegetation classes
and a standard confusion in this scene ($8{,}347$ and $1{,}322$ admissible training
pixels). The readiness scan reproduces the tissue structure: with standardized inputs a
random $12$-dimensional encoder reads the contrast at $0.850$ against a discriminant of
$0.933$; after PCA whitening the random encoder reads $0.667$ against $0.955$
(Table~\ref{tab:pavia_scan}), so the same two regimes (high and low initial probe
accessibility) apply, with a smaller asymmetry than on tissue. Calibration matched the
constructed contextual discriminant to the centre-only discriminant in each regime
($\tau=1.89$ standardized; $1.66$ whitened at $\gamma=30$ and at $\gamma=10$). Sizes
are scaled to the pair: up to $2{,}400$ balanced training patches, $400$ probe-fit
pixels, evaluation on all validation and test pixels of the pair. Arms, widths, rate
multipliers, seeds, thresholds, conditions and the retraining protocol are those of
Appendix~\ref{app:exp3}.

\subsection{Experiment 4 on Pavia: natural context}
All nine classes, macro-F1, with the Stage A gate and the Stage B arms of
Appendix~\ref{app:exp4}. Stage A gives the opposite of the tissue result
(Table~\ref{tab:exp4_pixel_paviau}): the best per-pixel classifier gains only $0.020$
macro-F1 on validation and $0.009$ on test from the components beyond the sixteenth
(MLP $0.891\to0.911$ and $0.876\to0.885$; discriminant $0.767\to0.781$ and
$0.742\to0.759$), so the registered gate fails and sixteen principal components carry
essentially all per-pixel information available to these classifiers on this scene,
whose bands are strongly correlated. By the registered rule the compression arms are
dropped here and Stage B runs the usage arms and the no-bottleneck arms.

% generated by paper/figures_src/make_tables_pavia.py
\begin{table}[htbp]\centering\scriptsize
\caption{Pavia University readiness scan (tile-disjoint halves of the training split): centre-only discriminant (oracle) and the discriminant on a random $K$-dimensional linear encoder's output (mean [min, max] over three encoders), standardized and PCA-whitened inputs; $|\cos|$ is the cosine between the class-mean contrast and the first principal direction.}\label{tab:pavia_scan}
\begin{tabular}{llcccccc}\toprule
pair & preprocessing & oracle & $K=2$ & $K=4$ & $K=12$ & $K=32$ & $|\cos|$ \\ \midrule
Meadows vs Trees & standardized & 0.933 & 0.666 [0.62, 0.71] & 0.731 [0.68, 0.78] & 0.850 [0.84, 0.88] & 0.921 [0.91, 0.92] & 0.91 \\
Meadows vs Trees & whitened & 0.955 & 0.576 [0.57, 0.58] & 0.580 [0.57, 0.59] & 0.667 [0.62, 0.70] & 0.815 [0.77, 0.87] & 0.91 \\
Gravel vs Bricks & standardized & 0.853 & 0.630 [0.61, 0.65] & 0.638 [0.62, 0.66] & 0.780 [0.73, 0.82] & 0.837 [0.82, 0.85] & 0.88 \\
Gravel vs Bricks & whitened & 0.833 & 0.519 [0.50, 0.55] & 0.559 [0.51, 0.60] & 0.612 [0.59, 0.64] & 0.676 [0.66, 0.69] & 0.88 \\
Asphalt vs Bricks & standardized & 0.952 & 0.846 [0.74, 0.91] & 0.898 [0.89, 0.91] & 0.922 [0.92, 0.93] & 0.944 [0.94, 0.95] & 0.99 \\
Asphalt vs Bricks & whitened & 0.945 & 0.544 [0.50, 0.57] & 0.566 [0.55, 0.57] & 0.650 [0.59, 0.70] & 0.775 [0.75, 0.82] & 0.99 \\
Meadows vs BareSoil & standardized & 0.910 & 0.750 [0.74, 0.75] & 0.779 [0.77, 0.79] & 0.837 [0.82, 0.85] & 0.908 [0.90, 0.91] & 0.99 \\
Meadows vs BareSoil & whitened & 0.925 & 0.605 [0.57, 0.66] & 0.662 [0.64, 0.69] & 0.729 [0.71, 0.75] & 0.797 [0.78, 0.82] & 0.99 \\
\bottomrule\end{tabular}\end{table}
\IfFileExists{sections_iclr_v2/appendix_exp3_paviau_tables.tex}{\begin{table}[htbp]\centering\scriptsize
\caption{Experiment 3, unready (whitened, $\gamma=30$) regime, at $L^\ast=0.30$: mean$\pm$sd over seeds (number of seeds). Probe gain = discriminant accuracy on the encoder output of validation centres minus its initial value; accuracies on validation (val) and test (test) patients under the paired conditions.}\label{tab:exp3_unready_paviau}
\begin{tabular}{llrcccccc}\toprule
arm & $M$ / $\kappa$ & step & probe gain & rev.\ (val) & ctx-rand.\ (val) & spec.-only (val) & rev.\ (test) & ctx-rand.\ (test) \\ \midrule
\texttt{sp} & 8 & 85$\pm$71 & -0.002$\pm$0.002 & 0.073$\pm$0.015 & 0.466$\pm$0.034 & 0.355$\pm$0.122 & 0.068$\pm$0.016 & 0.463$\pm$0.033 (3) \\
\texttt{sp} & 32 & 42$\pm$5 & -0.001$\pm$0.002 & 0.071$\pm$0.022 & 0.476$\pm$0.024 & 0.386$\pm$0.094 & 0.066$\pm$0.024 & 0.469$\pm$0.026 (3) \\
\texttt{sp} & 128 & 31$\pm$3 & 0.001$\pm$0.003 & 0.118$\pm$0.025 & 0.523$\pm$0.019 & 0.535$\pm$0.106 & 0.126$\pm$0.033 & 0.526$\pm$0.021 (3) \\
\texttt{sp} & 512 & 15$\pm$2 & -0.001$\pm$0.001 & 0.112$\pm$0.014 & 0.513$\pm$0.018 & 0.482$\pm$0.065 & 0.109$\pm$0.010 & 0.508$\pm$0.010 (3) \\
\midrule
\texttt{headlr} & 32 / 0.00390625 & 250$\pm$23 & 0.002$\pm$0.020 & 0.080$\pm$0.016 & 0.472$\pm$0.029 & 0.373$\pm$0.129 & 0.072$\pm$0.025 & 0.466$\pm$0.042 (3) \\
\texttt{headlr} & 32 / 0.0625 & 188$\pm$10 & 0.000$\pm$0.013 & 0.079$\pm$0.017 & 0.476$\pm$0.026 & 0.375$\pm$0.119 & 0.069$\pm$0.025 & 0.467$\pm$0.038 (3) \\
\texttt{headlr} & 32 & 42$\pm$5 & -0.001$\pm$0.002 & 0.071$\pm$0.022 & 0.476$\pm$0.024 & 0.386$\pm$0.094 & 0.066$\pm$0.024 & 0.469$\pm$0.026 (3) \\
\midrule
\texttt{ctxfree} & 8 & 10187$\pm$582 & 0.239$\pm$0.043 & 0.877$\pm$0.029 & 0.873$\pm$0.018 & 0.622$\pm$0.100 & 0.894$\pm$0.030 & 0.895$\pm$0.014 (3) \\
\texttt{ctxfree} & 128 & 9142$\pm$764 & 0.223$\pm$0.033 & 0.851$\pm$0.015 & 0.857$\pm$0.017 & 0.745$\pm$0.021 & 0.872$\pm$0.010 & 0.874$\pm$0.018 (3) \\
\midrule
\texttt{frozen} & 32 & 53$\pm$8 & 0.000$\pm$0.000 & 0.072$\pm$0.023 & 0.476$\pm$0.023 & 0.389$\pm$0.082 & 0.067$\pm$0.024 & 0.469$\pm$0.025 (3) \\
\bottomrule
\end{tabular}\end{table}
\begin{table}[htbp]\centering\scriptsize
\caption{Experiment 3, unready (whitened, $\gamma=10$; amendment) regime, at $L^\ast=0.30$: mean$\pm$sd over seeds (number of seeds). Probe gain = discriminant accuracy on the encoder output of validation centres minus its initial value; accuracies on validation (val) and test (test) patients under the paired conditions.}\label{tab:exp3_unready10_paviau}
\begin{tabular}{llrcccccc}\toprule
arm & $M$ / $\kappa$ & step & probe gain & rev.\ (val) & ctx-rand.\ (val) & spec.-only (val) & rev.\ (test) & ctx-rand.\ (test) \\ \midrule
\texttt{sp} & 8 & 416$\pm$244 & 0.008$\pm$0.006 & 0.056$\pm$0.008 & 0.479$\pm$0.019 & 0.365$\pm$0.106 & 0.053$\pm$0.010 & 0.479$\pm$0.018 (3) \\
\texttt{sp} & 32 & 283$\pm$28 & 0.004$\pm$0.005 & 0.054$\pm$0.012 & 0.486$\pm$0.011 & 0.391$\pm$0.064 & 0.052$\pm$0.016 & 0.484$\pm$0.015 (3) \\
\texttt{sp} & 128 & 220$\pm$20 & 0.008$\pm$0.009 & 0.071$\pm$0.013 & 0.506$\pm$0.011 & 0.514$\pm$0.081 & 0.076$\pm$0.014 & 0.506$\pm$0.009 (3) \\
\texttt{sp} & 512 & 103$\pm$16 & 0.001$\pm$0.003 & 0.070$\pm$0.007 & 0.500$\pm$0.009 & 0.481$\pm$0.050 & 0.070$\pm$0.003 & 0.500$\pm$0.006 (3) \\
\midrule
\texttt{headlr} & 32 / 0.00390625 & 1901$\pm$96 & 0.029$\pm$0.046 & 0.059$\pm$0.006 & 0.479$\pm$0.015 & 0.379$\pm$0.133 & 0.056$\pm$0.014 & 0.475$\pm$0.026 (3) \\
\texttt{headlr} & 32 / 0.0625 & 1303$\pm$79 & 0.021$\pm$0.035 & 0.056$\pm$0.009 & 0.481$\pm$0.013 & 0.381$\pm$0.110 & 0.054$\pm$0.013 & 0.478$\pm$0.023 (3) \\
\texttt{headlr} & 32 & 283$\pm$28 & 0.004$\pm$0.005 & 0.054$\pm$0.012 & 0.486$\pm$0.011 & 0.391$\pm$0.064 & 0.052$\pm$0.016 & 0.484$\pm$0.015 (3) \\
\midrule
\texttt{ctxfree} & 8 & 10067$\pm$1374 & 0.240$\pm$0.044 & 0.842$\pm$0.027 & 0.867$\pm$0.019 & 0.687$\pm$0.053 & 0.867$\pm$0.034 & 0.892$\pm$0.022 (3) \\
\texttt{ctxfree} & 128 & 8923$\pm$880 & 0.231$\pm$0.037 & 0.899$\pm$0.005 & 0.900$\pm$0.011 & 0.792$\pm$0.016 & 0.920$\pm$0.003 & 0.922$\pm$0.004 (3) \\
\midrule
\texttt{frozen} & 32 & 368$\pm$42 & 0.000$\pm$0.000 & 0.054$\pm$0.011 & 0.484$\pm$0.009 & 0.389$\pm$0.050 & 0.053$\pm$0.014 & 0.485$\pm$0.013 (3) \\
\bottomrule
\end{tabular}\end{table}
\begin{table}[htbp]\centering\scriptsize
\caption{Experiment 3, ready (standardized) regime, at $L^\ast=0.30$: mean$\pm$sd over seeds (number of seeds). Probe gain = discriminant accuracy on the encoder output of validation centres minus its initial value; accuracies on validation (val) and test (test) patients under the paired conditions.}\label{tab:exp3_ready_paviau}
\begin{tabular}{llrcccccc}\toprule
arm & $M$ / $\kappa$ & step & probe gain & rev.\ (val) & ctx-rand.\ (val) & spec.-only (val) & rev.\ (test) & ctx-rand.\ (test) \\ \midrule
\texttt{sp} & 8 & 1219$\pm$441 & 0.002$\pm$0.002 & 0.227$\pm$0.050 & 0.530$\pm$0.032 & 0.605$\pm$0.059 & 0.225$\pm$0.049 & 0.533$\pm$0.028 (3) \\
\texttt{sp} & 32 & 1066$\pm$262 & -0.001$\pm$0.004 & 0.230$\pm$0.021 & 0.544$\pm$0.010 & 0.634$\pm$0.036 & 0.227$\pm$0.027 & 0.547$\pm$0.019 (3) \\
\texttt{sp} & 128 & 842$\pm$77 & 0.000$\pm$0.003 & 0.264$\pm$0.035 & 0.574$\pm$0.028 & 0.733$\pm$0.031 & 0.279$\pm$0.046 & 0.583$\pm$0.029 (3) \\
\texttt{sp} & 512 & 431$\pm$65 & 0.000$\pm$0.002 & 0.252$\pm$0.017 & 0.568$\pm$0.017 & 0.719$\pm$0.036 & 0.255$\pm$0.009 & 0.569$\pm$0.015 (3) \\
\midrule
\texttt{headlr} & 32 / 0.00390625 & 13524$\pm$4802 & 0.003$\pm$0.004 & 0.189$\pm$0.028 & 0.512$\pm$0.041 & 0.520$\pm$0.136 & 0.181$\pm$0.043 & 0.513$\pm$0.046 (3) \\
\texttt{headlr} & 32 / 0.0625 & 5910$\pm$2023 & 0.001$\pm$0.004 & 0.227$\pm$0.030 & 0.537$\pm$0.030 & 0.591$\pm$0.080 & 0.219$\pm$0.038 & 0.536$\pm$0.029 (3) \\
\texttt{headlr} & 32 & 1066$\pm$262 & -0.001$\pm$0.004 & 0.230$\pm$0.021 & 0.544$\pm$0.010 & 0.634$\pm$0.036 & 0.227$\pm$0.027 & 0.547$\pm$0.019 (3) \\
\midrule
\texttt{ctxfree} & 8 & 7465$\pm$1642 & 0.004$\pm$0.002 & 0.853$\pm$0.033 & 0.855$\pm$0.027 & 0.834$\pm$0.013 & 0.941$\pm$0.028 & 0.942$\pm$0.022 (3) \\
\texttt{ctxfree} & 128 & 4781$\pm$463 & 0.003$\pm$0.002 & 0.875$\pm$0.002 & 0.877$\pm$0.006 & 0.869$\pm$0.009 & 0.964$\pm$0.002 & 0.962$\pm$0.004 (3) \\
\midrule
\texttt{frozen} & 32 & 1418$\pm$232 & 0.000$\pm$0.000 & 0.217$\pm$0.016 & 0.542$\pm$0.005 & 0.637$\pm$0.027 & 0.213$\pm$0.016 & 0.542$\pm$0.011 (3) \\
\bottomrule
\end{tabular}\end{table}
\begin{table}[htbp]\centering\scriptsize
\caption{Recovery at $M=32$ (real spectra): a fresh patch head with paired initialization trained for 20,000 steps on newly sampled context-random patches with the encoder frozen at its $L^\ast$ checkpoint (\texttt{init} = random initial encoder); validation and test patients under the paired conditions; original = the run's own classifier at $L^\ast$ (validation).}\label{tab:exp3_recovery_paviau}
\begin{tabular}{lllcccccc}\toprule regime & seed & encoder & probe (val) & orig.\ rev.\ (val) & retr.\ iid (val) & retr.\ rev.\ (val) & retr.\ ctx-rand.\ (val) & retr.\ rev.\ (test) \\ \midrule
ready & 0 & \texttt{init} & 0.915 & -- & 0.893 & 0.894 & 0.892 & 0.972 \\
ready & 0 & \texttt{kappa1} & 0.917 & 0.208 & 0.898 & 0.897 & 0.895 & 0.972 \\
ready & 0 & \texttt{kappa256} & 0.916 & 0.210 & 0.912 & 0.911 & 0.912 & 0.971 \\
ready & 1 & \texttt{init} & 0.923 & -- & 0.892 & 0.887 & 0.890 & 0.967 \\
ready & 1 & \texttt{kappa1} & 0.918 & 0.249 & 0.891 & 0.886 & 0.889 & 0.967 \\
ready & 1 & \texttt{kappa256} & 0.923 & 0.200 & 0.907 & 0.906 & 0.906 & 0.966 \\
ready & 2 & \texttt{init} & 0.909 & -- & 0.875 & 0.875 & 0.875 & 0.965 \\
ready & 2 & \texttt{kappa1} & 0.909 & 0.233 & 0.880 & 0.876 & 0.880 & 0.965 \\
ready & 2 & \texttt{kappa256} & 0.917 & 0.158 & 0.899 & 0.904 & 0.899 & 0.969 \\
unready & 0 & \texttt{init} & 0.664 & -- & 0.639 & 0.634 & 0.636 & 0.621 \\
unready & 0 & \texttt{kappa1} & 0.665 & 0.095 & 0.640 & 0.638 & 0.635 & 0.626 \\
unready & 0 & \texttt{kappa256} & 0.688 & 0.098 & 0.626 & 0.650 & 0.635 & 0.646 \\
unready & 1 & \texttt{init} & 0.703 & -- & 0.615 & 0.668 & 0.638 & 0.652 \\
unready & 1 & \texttt{kappa1} & 0.701 & 0.051 & 0.611 & 0.667 & 0.639 & 0.650 \\
unready & 1 & \texttt{kappa256} & 0.692 & 0.068 & 0.588 & 0.646 & 0.610 & 0.639 \\
unready & 2 & \texttt{init} & 0.755 & -- & 0.641 & 0.679 & 0.652 & 0.688 \\
unready & 2 & \texttt{kappa1} & 0.752 & 0.068 & 0.645 & 0.670 & 0.652 & 0.684 \\
unready & 2 & \texttt{kappa256} & 0.746 & 0.074 & 0.612 & 0.652 & 0.629 & 0.651 \\
unready10 & 0 & \texttt{init} & 0.664 & -- & 0.653 & 0.646 & 0.650 & 0.644 \\
unready10 & 0 & \texttt{kappa1} & 0.673 & 0.065 & 0.658 & 0.660 & 0.659 & 0.639 \\
unready10 & 0 & \texttt{kappa256} & 0.746 & 0.065 & 0.706 & 0.720 & 0.706 & 0.718 \\
unready10 & 1 & \texttt{init} & 0.703 & -- & 0.649 & 0.708 & 0.678 & 0.665 \\
unready10 & 1 & \texttt{kappa1} & 0.704 & 0.041 & 0.644 & 0.701 & 0.676 & 0.666 \\
unready10 & 1 & \texttt{kappa256} & 0.702 & 0.052 & 0.624 & 0.669 & 0.647 & 0.680 \\
unready10 & 2 & \texttt{init} & 0.755 & -- & 0.663 & 0.708 & 0.681 & 0.709 \\
unready10 & 2 & \texttt{kappa1} & 0.755 & 0.057 & 0.665 & 0.701 & 0.680 & 0.706 \\
unready10 & 2 & \texttt{kappa256} & 0.761 & 0.059 & 0.666 & 0.702 & 0.680 & 0.712 \\
\bottomrule\end{tabular}\end{table}
\paragraph{Runs that did not reach $L^\ast$.}
Every run reached $L^\ast$.
}{}
\IfFileExists{sections_iclr_v2/appendix_exp4_paviau_tables.tex}{% generated by paper/figures_src/make_tables_exp4.py
\begin{table}[htbp]\centering\scriptsize
\caption{Experiment 4 (Pavia University), Stage A: four-class macro-F1 of per-pixel classifiers on the centre spectra against the number of principal components $k$ (PCA on standardized training-split centres). Shrinkage Fisher discriminant (LDA), balanced logistic regression (LR) and a 256-unit ReLU MLP with a fixed schedule (mean$\pm$sd over three seeds), on validation and test patients.}\label{tab:exp4_pixel_paviau}
\begin{tabular}{rcccccc}\toprule
$k$ & LDA (val) & LR (val) & MLP (val) & LDA (test) & LR (test) & MLP (test) \\ \midrule
2 & 0.549 & 0.572 & 0.594$\pm$0.006 & 0.719 & 0.723 & 0.731$\pm$0.003 \\
4 & 0.639 & 0.662 & 0.722$\pm$0.004 & 0.676 & 0.733 & 0.776$\pm$0.008 \\
8 & 0.663 & 0.741 & 0.790$\pm$0.005 & 0.688 & 0.810 & 0.826$\pm$0.006 \\
16 & 0.767 & 0.816 & 0.891$\pm$0.004 & 0.742 & 0.831 & 0.876$\pm$0.003 \\
23 & 0.783 & 0.825 & 0.902$\pm$0.004 & 0.757 & 0.838 & 0.882$\pm$0.004 \\
32 & 0.785 & 0.828 & 0.905$\pm$0.004 & 0.770 & 0.839 & 0.882$\pm$0.005 \\
64 & 0.784 & 0.829 & 0.907$\pm$0.005 & 0.762 & 0.840 & 0.882$\pm$0.003 \\
103 & 0.781 & 0.830 & 0.911$\pm$0.005 & 0.759 & 0.841 & 0.885$\pm$0.006 \\
\bottomrule\end{tabular}\end{table}
}{}

